\paragraph{La normalización de alfa y su configuración firmada.}
\label{inc-ampl-clausura}
La coordenada de $\alpha$ se determina en la vía posicional mediante la
composición de los bloques regionales y del registro dodecafásico. Antes
de la división posicional, la combinación orientada tiene componentes
$Z_j=P_j+E_j-\Phi_j-K_j^{\mathrm{per}}$ y el balance con cociente entrante
es $s_j=Z_j+c_{j+1}$. La normalización satisface
\begin{equation}
 A_j=s_j-1000c_j=Z_j+c_{j+1}-1000c_j,\qquad 0\leq A_j<1000.
 \label{inc-ampl-normalizacion}
\end{equation}
Esta identidad conserva simultáneamente el bloque normalizado y la
contribución del cociente transportado. El límite de los prefijos
compatibles determina la coordenada escalar de la clausura. La
configuración firmada conserva, por su parte, qué posiciones de la
combinación anterior a la normalización presentan defecto negativo.
En la ventana inicial, esa configuración es el vector $Z_0$ de
\eqref{exc:precarry}, y su soporte negativo es $H^-_\alpha$.

Así quedan definidos dos contenidos de una misma operación. La
normalización posicional determina una coordenada real mediante sus
restos y cocientes compatibles; el soporte negativo determina una
configuración de posiciones en el ciclo de doce componentes. El signo se
refiere al balance de registros anterior a la normalización. Su lectura
incidencial permanece disponible porque el estado conserva ese balance
junto con el resultado normalizado. Esta articulación corresponde a la
vía posicional demostrada; la comparación con los operadores analíticos
de vacancias conserva el alcance preciso establecido en el capítulo de
$\alpha$.
